% =====================================================================
\chapter{C1: TA-offline authentication with per-member key delivery}\label{ch:c1}
% =====================================================================
\begin{keybox}[colframe=uavc,colback=uavc!5]
The TA works only \textbf{before} the mission. Each entity holds a \emph{certificateless} key $\sk=x+d$: the TA issues the partial key $d$, the entity adds
its own secret $x$ that never leaves it. In the field the TUAV checks all join requests with the TA's public key alone, in one weighted batch, and
sends each accepted UAV the group key \emph{sealed under its own fresh Diffie--Hellman key}. No satellite, no pairing, no public polynomial.
\end{keybox}

\section{Design goals}
\begin{center}
\small
\begin{tabularx}{\linewidth}{@{}l X l@{}}
\toprule
& \textbf{Goal} & \textbf{Answers}\\
\midrule
G1 & the authority that issues credentials cannot sign for anybody & F1\\
G2 & nothing published depends on the members' secrets; removed members and newcomers learn nothing & F2\\
G3 & every key delivery uses fresh, erased secrets & F3\\
G4 & authentication needs only public parameters, so the TA can stay offline during the mission & design\\
G5 & no pairings; many requests checked at once, and one bad request cannot block the others & design\\
\bottomrule
\end{tabularx}
\end{center}

\begin{figure}[H]
\centering
\begin{tikzpicture}[font=\sffamily\small]
  \node[draw=black,dashed,rounded corners=5pt,fill=black!4,minimum width=5.2cm,minimum height=4.4cm] (pre) at (0,0) {};
  \node[font=\sffamily\bfseries] at (0,1.85) {Before the mission (TA online)};
  \pic[scale=1.3] at (0,0.75) {ta};
  \node[font=\scriptsize\sffamily,align=center] at (0,-0.75) {Setup: master key $s$, tracing key $t$\\ Registration: one credential per\\ epoch for every UAV and TUAV};
  \node[draw=black,dashed,rounded corners=5pt,fill=white,minimum width=8.6cm,minimum height=4.4cm] (fld) at (9.0,0) {};
  \node[font=\sffamily\bfseries] at (9.0,1.85) {During the mission (TA offline)};
  \pic[scale=1.3] at (5.9,0.45) {drone};
  \node[font=\scriptsize\sffamily\bfseries] at (5.9,-0.35) {TUAV};
  \foreach \x/\y in {10.6/1.0,11.8/0.45,10.9/-0.25,12.2/-0.7}{\pic[scale=0.8] at (\x,\y) {drone};}
  \node[font=\scriptsize\sffamily\bfseries] at (11.4,-1.25) {UAVs};
  \draw[{Stealth[length=2.6mm]}-{Stealth[length=2.6mm]},very thick] (6.75,0.45) -- node[above,font=\scriptsize\sffamily,align=center]{join UA1--UA5:\\ one local round trip} (10.0,0.45);
  \node[font=\scriptsize\sffamily,align=center] at (9.0,-1.8) {verify with $\Ppub$ only \;$\cdot$\; weighted batch check \;$\cdot$\; one sealed key per member};
  \draw[-{Stealth[length=3mm]},very thick] (pre.east) -- node[above,font=\scriptsize\sffamily]{credentials} (fld.west);
\end{tikzpicture}
\caption{Overview of contribution C1. Left: before the mission, the TA sets up the system and issues one certificateless credential to every UAV and TUAV.
Right: during the mission the TA is offline; a UAV joins its TUAV in one local exchange (steps UA1--UA5), the TUAV verifies requests using only the
TA's public key $\Ppub$, and each member receives the group key sealed under its own key.}
\label{fig:c1big}
\end{figure}

\section{Keys: certificateless, so the TA cannot sign}
\begin{stepbox}{deep}{ST and RG1--RG4 \;(before the mission)}
\steplabel{deep}{ST.} The TA picks a master key $s$ and a tracing key $t$, and publishes $\Ppub=sP$ and $\Tpub=tP$.\\
\steplabel{deep}{RG1.} The entity picks its own secret $x$ and sends $X=xP$ with a proof of possession. \textbf{$x$ never leaves the entity.}\\
\steplabel{deep}{RG2.} The TA picks $\rho,y$ and computes a pseudonym and a partial key
\begin{gather*}
\PID=(\PID^1,\PID^2)=\big(\rho P,\ \RID\oplus\Hz(\rho\Tpub\cat\ep)\big),\\
Y=yP,\qquad h=\Ho(\PID\cat\ep\cat X\cat Y),\qquad d=y+s\,h\bmod q ,
\end{gather*}
where the label $\ep$ encodes the validity epoch and the role ($\ep=0$ is reserved for TUAVs).\\
\steplabel{deep}{RG3.} The TA sends $(\PID,\ep,Y,d)$ to the entity over the registration channel.\\
\steplabel{deep}{RG4.} The entity checks $dP\overset{?}{=}Y+h\Ppub$ and keeps $\sk=x+d$, \;$\pk=X+Y+h\Ppub$.
\end{stepbox}

\begin{lemma}[Key consistency]\label{lem:keycons}
$\sk\cdot P=\pk$, and anyone can compute $\pk$ from the public tuple $(\PID,\ep,X,Y)$ and $\Ppub$.
\end{lemma}
\begin{proof}
$\sk\cdot P=xP+yP+s\,hP=X+Y+h\Ppub=\pk$, and $h=\Ho(\PID\cat\ep\cat X\cat Y)$ is public.
\end{proof}
\noindent The pair $(Y,d)$ is a Schnorr signature of the TA on $(\PID,\ep,X)$; this is the Galindo--Garcia identity-based construction~\cite{gg09}, made
certificateless by adding the user's own $x$~\cite{alriyami}. Only the tracing authority, which knows $t$, can compute $\rho\Tpub=t\,\PID^1$ and open
$\PID^2$; a fresh $\rho$ per epoch makes pseudonyms of different epochs unlinkable.

\section{Joining a TUAV: UA1--UA5}
\begin{stepbox}{uavc}{UA1--UA5 \;(in the field, TA offline)}
\steplabel{hoc}{UA1.} TUAV $\Ta$ picks $e_a$, sets $E_a=e_aP$ and broadcasts a beacon $\langle\ID_a,X_a,Y_a,E_a,\ts_a\rangle$ signed with its own key.\\
\steplabel{uavc}{UA2.} UAV $\U_i$ verifies the beacon, picks $e_i,r_i$ and sends $\langle\PID_i,\ep,X_i,Y_i,E_i,\ts_i,R_i,v_i\rangle$, where
\begin{gather*}
E_i=e_iP,\qquad R_i=r_iP,\qquad m_i=(\texttt{join}\cat\ID_a\cat E_a\cat E_i),\\
c_i=\Ht(m_i\cat\PID_i\cat\ep\cat X_i\cat Y_i\cat R_i\cat\ts_i),\qquad v_i=r_i+c_i\,\sk_i .
\end{gather*}
\steplabel{hoc}{UA3.} The TUAV draws random 64-bit weights $\lambda_i$ and checks all $n$ requests at once:
\begin{equation}\label{eq:bv}
\Big(\sum_{i}\lambda_iv_i\Big)P\overset{?}{=}\sum_i\lambda_iR_i+\sum_i(\lambda_ic_i)(X_i+Y_i)+\Big(\sum_i\lambda_ic_ih_i\Big)\Ppub .
\end{equation}
If the check fails, it verifies each request alone, accepts the valid ones and lists the invalid pseudonyms (bounded fallback).\\
\steplabel{hoc}{UA4.} For every accepted $\U_i$: \;$K_i=\KDF(e_aE_i\cat\PID_i\cat\ID_a\cat\ts_i)$ and
$C_i=\Enc\big(K_i,\ \GK\cat\nu;\ \Hd(\PID_i)\cat\nu\cat\ts_3\big)$. All entries go out in \textbf{one} broadcast $\langle\ID_a,\ts_3,\nu,\{\Hd(\PID_i),C_i\}_i\rangle$.\\
\steplabel{uavc}{UA5.} $\U_i$ computes $K_i=\KDF(e_iE_a\cat\PID_i\cat\ID_a\cat\ts_i)$, opens $C_i$ (which also proves the TUAV knows $e_a$), stores
$(\GK,\nu)$ and \textbf{erases} $e_i$. The TUAV erases $e_a$ after the round.
\end{stepbox}

\begin{figure}[H]
\centering
\resizebox{\linewidth}{!}{%
\begin{tikzpicture}[font=\sffamily\small]
  \node[ent=hoc]  (t) at (0,0)  {TUAV $\Ta$};
  \node[ent=uavc] (u) at (11,0) {UAV $\U_i$ \;($\times n$)};
  \node[ent=deep,dashed,minimum width=2cm] (ta) at (5.5,0) {TA (offline)};
  \draw[life] (t.south) -- ++(0,-9.6);  \draw[life] (u.south) -- ++(0,-9.6);
  \draw[black!15,line width=4pt] (ta.south) -- ++(0,-9.6);
  \node[font=\scriptsize\sffamily,text=black!45,rotate=90] at (5.5,-5.6) {never contacted};
  \node[act=hoc,anchor=north] at (0,-0.9) {UA1: pick $e_a$, $E_a=e_aP$;\\ sign $(\ID_a,E_a,\ts_a)$};
  \draw[msg=hoc] (0,-2.0) -- node[lbl,above]{UA1 \;signed beacon $\langle \ID_a,X_a,Y_a,E_a,\ts_a,R_a,v_a\rangle$} (11,-2.0);
  \node[act=uavc,anchor=north] at (11,-2.3) {verify beacon; pick $e_i,r_i$\\ $E_i=e_iP$, \;$R_i=r_iP$, \;$v_i=r_i+c_i\sk_i$};
  \draw[msg=uavc] (11,-3.9) -- node[lbl,above]{UA2 \;$\langle \PID_i,\ep,X_i,Y_i,E_i,\ts_i,R_i,v_i\rangle$} (0,-3.9);
  \node[act=hoc,anchor=north] at (0,-4.2) {UA3: weighted batch check (\ref{eq:bv});\\ on failure $\to$ individual fallback};
  \node[act=hoc,anchor=north] at (0,-5.4) {UA4: $K_i=\KDF(e_aE_i\cat\dots)$\\ $C_i=\Enc(K_i,\GK\cat\nu)$};
  \draw[msg=hoc] (0,-6.8) -- node[lbl,above]{UA4 \;$\langle \ID_a,\ts_3,\nu,\{\Hd(\PID_i),C_i\}\rangle$ \;(one broadcast)} (11,-6.8);
  \node[act=uavc,anchor=north] at (11,-7.1) {UA5: $K_i=\KDF(e_iE_a\cat\dots)$; open $C_i$\\ $\Rightarrow$ TUAV authentic; erase $e_i$};
  \node[draw=good,fill=good!10,rounded corners,font=\scriptsize\sffamily\bfseries,text=good] at (5.5,-9.0) {mutual authentication + shared $\GK$ in one local round trip};
\end{tikzpicture}}
\caption{Message flow of C1 between the TUAV and $n$ UAVs. Time runs from top to bottom; boxes show local computation. The TA (centre) is never contacted. Compared with Figure~\ref{fig:baseflow}, there is no satellite link, no pairing and no CRT polynomial.}
\label{fig:c1flow}
\end{figure}

\section{Rekeying on every change}
\begin{stepbox}{hoc}{RK1--RK2 \;(join and leave)}
\steplabel{hoc}{RK1 (leave or eviction).} The TUAV draws a \emph{fresh} $\GK'$, sets $\nu\leftarrow\nu+1$ and sends $C_i'=\Enc(K_i,\GK'\cat\nu)$ only to
the members that stay.\\
\steplabel{hoc}{RK2 (join).} After admitting newcomers (UA1--UA5), the TUAV immediately rekeys the whole cell, so a newcomer never receives a key that was
already in use (fix A6). Members accept only increasing versions $\nu$ (fix A5).
\end{stepbox}
\begin{remark}[Large groups]
Each rekey costs one entry per member. For large cells the TUAV can arrange keys in a logical key hierarchy (LKH)~\cite{wgl98}, which reduces a rekey to
$O(\log n)$ entries; nothing else changes.
\end{remark}

\section{Correctness}
\begin{lemma}[Signature and batch correctness]\label{lem:c1corr}
(i) An honest request satisfies $v_iP=R_i+c_i\pk_i$. (ii) A batch of honest requests satisfies Eq.~\eqref{eq:bv} for every choice of weights.
(iii) Both sides of UA4/UA5 derive the same $K_i$.
\end{lemma}
\begin{proof}
(i) $v_iP=r_iP+c_i\sk_iP=R_i+c_i\pk_i$ by Lemma~\ref{lem:keycons}. (ii) Multiply (i) by $\lambda_i$, sum over $i$, and expand
$\pk_i=X_i+Y_i+h_i\Ppub$. (iii) $e_aE_i=e_ae_iP=e_iE_a$, and the other KDF inputs are identical.
\end{proof}


\section{Security}
\begin{theorem}[Unforgeability]\label{thm:unforge}
In the random-oracle model, if the ECDLP is hard in $\G$, no PPT adversary can output a valid signature on a new message under the credential
$(\PID_i,\ep,X_i,Y_i)$ of an honest entity -- neither \textbf{(Type I)} an outsider who does not know $s$ but may register its own keys and replace other
public keys, nor \textbf{(Type II)} a malicious TA who knows $s$ and every partial key $d_i$.
\end{theorem}
\begin{proof}[Proof sketch]
\emph{Type II.} Given an ECDLP challenge $X^*$, the reduction plays the TA honestly (it knows $s$ and chooses $y$) and sets $X_i=X^*$ for one guessed
honest entity. It answers signing queries without $\sk_i$ by the standard Schnorr simulation: pick $v,c$, set $R=vP-c\,\pk_i$ and program $\Ht$. A forgery
is rewound with the General Forking Lemma (\ref{lem:fork}) on $\Ht$, which yields $\sk_i=x_i+d_i$; since $d_i$ is known, $x_i=\sk_i-d_i$ solves the
challenge.\\
\emph{Type I.} If the forgery is under the honest $(X_i,Y_i)$, the same argument applies with the challenge planted in $X_i$. If the adversary uses keys
$(X',Y')$ of its own, the proof of possession of $X'$ gives the reduction $x'$, and forking on $\Ht$ yields $\sk'=x'+d'$, hence $d'$ with
$d'P=Y'+h'\Ppub$. Because $h'=\Ho(\PID\cat\ep\cat X'\cat Y')$ is fixed only after $Y'$, a second fork on $\Ho$ gives $d''$ with $d''P=Y'+h''\Ppub$, and
$s=(d'-d'')(h'-h'')^{-1}$ solves an ECDLP challenge planted in $\Ppub$. This double-forking structure is that of the Galindo--Garcia
scheme~\cite{gg09}; a careful analysis of its (non-tight) reduction is given in~\cite{ckk12}.
\end{proof}

\begin{theorem}[Batch soundness]\label{thm:batch}
If a batch contains at least one invalid request, Eq.~\eqref{eq:bv} holds with probability at most $2^{-64}$ over the TUAV's weights. With the fallback,
the TUAV accepts exactly the valid requests, so one forged request cannot block the honest ones.
\end{theorem}
\begin{proof}
Eq.~\eqref{eq:bv} is $\sum_i\lambda_i\Delta_i=0$ with $\Delta_i=v_iP-R_i-c_i\pk_i$; an invalid request has $\Delta_i\neq0$. Apply the Small-Exponent Test
(Lemma~\ref{lem:bgr}) with $\ell=64$. The fallback checks each request individually, which is exact.
\end{proof}

\begin{theorem}[No key escrow]\label{thm:noescrow}
The TA cannot produce a valid signature under an honest entity's credential $(X_i,Y_i)$ unless it solves the ECDLP. It can only issue a \emph{second}
credential $(X',Y')$ for the same pseudonym -- Girault trust level~2, compared with level~1 in the base scheme (Corollary~\ref{cor:escrow}).
\end{theorem}
\begin{proof}
This is the Type~II case of Theorem~\ref{thm:unforge}: the TA knows $d_i$ but not $x_i$, and $\sk_i=x_i+d_i$.
\end{proof}

\begin{theorem}[Group-key secrecy, revocation and backward secrecy]\label{thm:gksecrecy}
Assume CDH is hard in $\G$, $\KDF$ is a random oracle, $\Enc$ is a secure AEAD scheme and signatures are unforgeable (Theorem~\ref{thm:unforge}). Then the
group key $\GK_\nu$ of version $\nu$ is indistinguishable from random for (a) any active outsider, (b) any member removed before version $\nu$, and (c) any
member that joined after version $\nu$.
\end{theorem}
\begin{proof}[Proof sketch]
$\GK_\nu$ appears only inside the ciphertexts $C_i=\Enc(K_i,\GK_\nu\cat\nu)$ of the members of version $\nu$. Each $K_i$ is derived from $e_aE_i$, where
$E_a$ is signed by the TUAV (UA1) and $E_i$ is signed by the member inside $m_i$ (UA2), so an active adversary cannot substitute its own value without
forging a signature. Computing $e_aE_i$ from $E_a$ and $E_i$ is CDH; in the random-oracle model $K_i$ is therefore uniform to everyone except the two
endpoints, and AEAD security makes $C_i$ independent of $\GK_\nu$. (a) An outsider holds no $K_i$. (b) A removed member gets no entry for version $\nu$
(RK1); its old keys open only older entries, and $\GK_\nu$ is drawn fresh and independently. (c) A newcomer is rekeyed on arrival (RK2), so it receives a
fresh key independent of the earlier ones, which were sealed under other members' keys. Unlike the base scheme, no published value depends on any member's
secret, so Theorems~\ref{thm:insider} and~\ref{thm:outsider} have nothing to work on.
\end{proof}

\begin{theorem}[Forward secrecy]\label{thm:fs}
If every long-term secret ($s$, all $x_i$ and $d_i$, the TUAV's key) is revealed after version $\nu$ has ended and the exponents $e_a,e_i$ have been
erased, $\GK_\nu$ remains indistinguishable from random under CDH.
\end{theorem}
\begin{proof}
Long-term keys are used only to sign; they do not enter $K_i$. Recovering $K_i$ requires $e_ae_iP$ from the recorded $E_a=e_aP$ and $E_i=e_iP$, which is
a CDH instance; the rest follows as in Theorem~\ref{thm:gksecrecy}.
\end{proof}

\section{How C1 answers the three flaws}
\begin{center}
\small
\begin{tabularx}{\linewidth}{@{}l X l@{}}
\toprule
\textbf{Flaw} & \textbf{What C1 does, mathematically} & \textbf{Result}\\
\midrule
\textcolor{bad}{\textbf{F1}} TA impersonation & $\sk=x+d$; the TA never sees $x$ & Thm.~\ref{thm:noescrow}\\
\textcolor{bad}{\textbf{F2}} CRT leakage & no public function of members' secrets; one sealed entry per member; fresh key on every join and leave & Thm.~\ref{thm:gksecrecy}\\
\textcolor{bad}{\textbf{F3}} no forward secrecy & keys come from ephemeral $e_a,e_i$, erased after use & Thm.~\ref{thm:fs}\\
\bottomrule
\end{tabularx}
\end{center}

\section{Cost, computed with the base paper's own constants}
The base paper (Section~VI-A) prices its operations with the MIRACL library: pairing $T_P=7.2232$\,ms, scalar multiplication in $\G_1$
$T_{SM\cdot P}=2.6784$\,ms, point addition in $\G_1$ $T_{PA\cdot P}=0.0342$\,ms, scalar multiplication on the ordinary curve $T_{SM\cdot E}=0.9018$\,ms,
small-scalar multiplication $T_{SSM\cdot E}=0.0475$\,ms, point addition $T_{PA\cdot E}=0.0107$\,ms and hashing $T_{SH}=0.0012$\,ms. C1 needs only the
ordinary-curve operations.

\begin{center}
\small
\begin{tabularx}{\linewidth}{@{}p{3.3cm} X X@{}}
\toprule
\textbf{Authenticating and keying $n$ UAVs} & \textbf{Base scheme} & \textbf{C1}\\
\midrule
Pairings & $3n$ (at the TA, GA2) & \textbf{0}\\
Satellite round trips & 1 per login round & \textbf{0}\\
TA work per UAV & $3T_P+3T_{SM\cdot P}+T_{PA\cdot P}\approx29.74$\,ms & \textbf{0} (offline)\\
TUAV work & $1.2524n+0.9054$\,ms (the paper's own GA formula) & $(2n{+}4)T_{SM\cdot E}+nT_{SSM\cdot E}+3nT_{PA\cdot E}+4nT_{SH}$
$\approx1.888n+3.607$\,ms\\
Total for $n=120$ & $\approx3720$\,ms + satellite trip & $\approx\mathbf{230}$\,\textbf{ms}\\
UAV work & 1.82\,ms for US (as reported) + GD5 + GD6 & $6T_{SM\cdot E}+3T_{PA\cdot E}+5T_{SH}\approx5.45$\,ms\\
Key-distribution traffic & $\approx10n^3$ bytes (Remark~\ref{rem:size}) & $\approx80n$ bytes (one broadcast)\\
\bottomrule
\end{tabularx}
\end{center}
The UAV in C1 does more work than the paper's US phase because it also authenticates the TUAV and runs a Diffie--Hellman exchange; that is the price of
mutual authentication and forward secrecy, and it stays a few milliseconds. On a modern 128-bit curve all scalar multiplications become somewhat more
expensive, but C1 still needs no pairing and no satellite.

\begin{figure}[H]
\centering
\begin{tikzpicture}
\begin{axis}[width=0.86\linewidth,height=6.4cm,ymode=log,xmin=10,xmax=200,ymin=0.5,ymax=1e6,
  xlabel={number of UAVs $n$},ylabel={key-distribution traffic (KB)},grid=major,grid style={black!8},
  tick label style={font=\small},label style={font=\small},legend style={font=\small,at={(0.02,0.97)},anchor=north west,draw=black!30}]
  \addplot[very thick,color=plotB,domain=10:200,samples=60] {20*x*(x*(x-1)/2+x+1)/1000};
  \addlegendentry{base: $n$ ACKs, each with all coefficients}
  \addplot[very thick,color=plotA,domain=10:200,samples=60] {80*x/1000};
  \addlegendentry{C1: one broadcast, $\approx80$ bytes per member}
  \addplot[only marks,mark=*,mark size=2.5pt,color=plotB] coordinates {(120,17426)};
  \addplot[only marks,mark=*,mark size=2.5pt,color=plotA] coordinates {(120,9.6)};
  \node[font=\small,text=plotB,anchor=south east] at (axis cs:118,20000) {17.4\,MB};
  \node[font=\small,text=plotA,anchor=south east] at (axis cs:118,11) {9.6\,KB};
\end{axis}
\end{tikzpicture}
\caption{Traffic needed to distribute one group key against the number of UAVs $n$ (logarithmic vertical axis), computed from the message formats. The base scheme sends all polynomial coefficients to each UAV, about $10n^3$ bytes in total; C1 sends one broadcast with about 80 bytes per member. The marked points are the values for $n=120$.}
\label{fig:traffic}
\end{figure}

\begin{limitbox}
\begin{itemize}[leftmargin=*]
\item The TA can still enrol a \emph{phantom} member under a new pseudonym; C1 removes its power to sign as an \emph{existing} entity (level~2, not 3).
\item The security reductions use the forking lemma and are not tight~\cite{ckk12}; this is standard for Schnorr-type schemes.
\item After C1 the TUAV knows every $K_i$ and the group key: \textbf{the trust has moved from the TA to the TUAV}. C2--C4 protect exactly this point.
\end{itemize}
\end{limitbox}
